\chapter{ДУГААР ТАНИХ СУДАЛГААНЫ ТОЙМ БА ОНОЛЫН ҮНДЭС}
\section{Асуудлын ангилал ба хоёр шатлалт хандлага}
Улсын дугаар таних асуудлыг уламжлалт байдлаар (1) дугаарын байрлалыг олох (plate detection), (2) дугаарын тэмдэгтүүдийг унших (plate recognition) гэсэн хоёр шатанд хуваадаг~\cite{laroca2018}. Эрт үеийн системүүд өнгө, ирмэг, морфологи дээр суурилсан дүрэмт аргаар дугаарыг олж, тэмдэгт тус бүрийг сегментчилээд ангилагчаар таньдаг байсан бол гүн сургалтын үед илрүүлэлтийг нэг шатлалт детектор (YOLO, SSD), уншилтыг сегментчлэлгүй дарааллын загвар (CRNN+CTC, attention-based) гүйцэтгэх болсон. Сегментчлэлгүй уншилт нь бохир, бүдэг, налуу дугаарт тэмдэгт тус бүрийг салгах шаардлагагүй тул илүү тэсвэртэй.

\section{Нэг шатлалт илрүүлэгч: YOLO гэр бүл}
YOLO (You Only Look Once) нь зургийг нэг удаа сүлжээгээр нэвтрүүлж объектын хайрцаг, ангийг шууд таамагладаг. Ultralytics-ийн YOLO11~\cite{yolo11} хувилбар нь anchor-free толгой, C3k2 ба C2PSA блок, DFL (Distribution Focal Loss) хайрцгийн регрессийг ашигладаг; хамгийн жижиг n хувилбар нь 2.6 сая параметртэй, 640~px оролтод GPU дээр хэдхэн миллисекундэд ажилладаг. Дугаар илрүүлэлтэд YOLO-г ашигласан ажлууд олон~\cite{laroca2018,morse2025}; гэхдээ COCO өгөгдлийн 80 ангид дугаар байхгүй тул бэлэн жин дугаарыг илрүүлэхгүй, дугаарын өгөгдлөөр сургах шаардлагатай.

\section{Дарааллын уншигч: CRNN ба CTC}
Shi нарын CRNN~\cite{shi2017crnn} нь конволюцийн сүлжээгээр зургаас баганын онцлогийн дараалал гаргаж, хоёр чиглэлт LSTM-ээр контекст нэмж, CTC (Connectionist Temporal Classification)~\cite{graves2006ctc} алдагдлаар тэмдэгтийн байрлалыг заахгүйгээр бүхэл дарааллыг сургадаг. CTC нь $T$ алхамтай гаралтын дарааллыг blank тэмдэгт ба давталтыг нэгтгэх $\mathcal{B}$ хувиргалтаар зорилтот текст рүү буулгадаг:
\begin{equation}
\mathcal{L}_{\mathrm{CTC}} = -\sum_{(x,y)}\log\sum_{\pi\in\mathcal{B}^{-1}(y)}\prod_{t=1}^{T}p(\pi_t\mid x).
\end{equation}
PaddleOCR~\cite{ppocr} хэрэгслийн уншигч нь мөн энэ CRNN+CTC зарчмаар бүтээгдсэн (MobileNetV3/ResNet суурь + BiLSTM + CTC) бөгөөд олон улсад хамгийн өргөн хэрэглэгддэг хөнгөн OCR шийдэл юм. Энэ ажилд PaddleOCR-ийн уншигчийн архитектурын зарчмыг PyTorch дээр дахин хэрэгжүүлж, монгол кирилл үсгийн 35 үсэг + 10 орон гэсэн өөрийн тэмдэгтийн олонлог дээр сургасан; учир нь PaddleOCR-ийн бэлэн кирилл загварт Ө, Ү үсэг болон дугаарын дүрэм тусгагдаагүй, Paddle фреймворк нь шинэ Blackwell GPU дээр найдвартай ажиллахгүй байсан.

\section{Синтетик өгөгдөл}
Jaderberg нарын~\cite{jaderberg2014} синтетик текст үүсгэх арга нь шошготой бодит зураг цөөн үед OCR загварыг урьдчилан сургах стандарт шийдэл болсон. Дугаарын хувьд рендерингт суурилсан синтетик дугаар~\cite{ruder2022} нь фонт, эмблем, өнгө, гэрэл, өнцөг, бүдэг зэргийг удирдаж олон мянган хувилбар үүсгэх боломж олгодог. Синтетик нь тэмдэгтийн хэлбэрийг заадаг ч бодит камерын шуугиан, бохирдол, шахалтын онцлогийг бүрэн дуурайж чадахгүй тул бодит өгөгдлөөр нарийн тохируулах (fine-tuning) шаардлагатай байдаг.

\section{Олон кадрын баталгаажуулалт}
Laroca нар~\cite{laroca2018} нэг машины бүх кадрын уншилтыг нэгтгэж байрлал бүрт хамгийн олон давтагдсан тэмдэгтийг сонгох олонхийн саналаар видео ALPR-ийн нарийвчлалыг мэдэгдэхүйц ахиулсан. Wang нар~\cite{wang2026mt} 33 клип дээр нэг кадрын 66.7 хувиас саналаар 81.8 хувь болгосон; гэхдээ бараг бүх кадрт буруу уншигдсан дугаарыг санал засахгүй гэдгийг тэмдэглэсэн. Pham нар~\cite{pham2023ctsm} тэмдэгтийн цаг хугацааны дарааллыг тулгах аргаар UFPR-ALPR дээр 96.7 хувьд хүрсэн. Эдгээр ажлууд итгэлийн жин ба дугаарын дүрмийг саналд оруулах нь нэмэлт давуу тал өгөхийг санал болгодог бөгөөд энэ ажилд яг тийм жинлэсэн саналыг хэрэгжүүлэв.

\section{Монгол дугаарын судалгаа ба өгөгдлийн сан}
Монгол дугаарт зориулсан хэвлэгдсэн судалгааны өгүүлэл, нийтийн benchmark олдсонгүй. Нээлттэй кодын хувьд gumbernator-ийн Mongolian-ALPR~\cite{gumbernator} репо нь TensorFlow Object Detection API-ийн SSD/Faster R-CNN илрүүлэгч (66--72 mAP гэж мэдээлсэн) ба тэмдэгт тус бүрийн Keras CNN ангилагчтай боловч жингийн холбоос алдагдсан; 68 зурагт дугаар ба тэмдэгт бүрийн хайрцгийн шошго хавсаргасан нь энэ ажилд ашиглагдсан. Tsundgerhen-ийн репо~\cite{tsund} YOLOv8-аар дугаар илрүүлдэг бөгөөд дотор нь MLUB тэмцээний 7394 бодит дугаарын тайралт текстийн хамт байсан. narangoo9-ийн систем~\cite{narangoo} OpenCV ба Tesseract (\texttt{mon.traineddata}) ашигладаг ч хэмжүүр мэдээлээгүй. Roboflow Universe дээрх ``mongolian plate'' өгөгдөл~\cite{roboflowmn} 984 зурагтай, зөвхөн хайрцгийн шошготой, CC BY 4.0 лицензтэй. Олон улсын өгөгдлийн сангуудыг Хүснэгт~\ref{tab:intl}-д харьцуулав.
\begin{table}[H]\centering\caption{Дугаар таних олон улсын гол өгөгдлийн сангууд ба монгол эх үүсвэрүүд}\label{tab:intl}\small
\resizebox{\textwidth}{!}{\begin{tabular}{lrll}\toprule
Өгөгдлийн сан & Зураг & Шошго & Улс / тэмдэглэл\\\midrule
CCPD~\cite{ccpd} & 300\,000+ & хайрцаг + текст & Хятад\\
UFPR-ALPR~\cite{ufpr} & 4\,500 & хайрцаг + текст, видео & Бразил, хүсэлтээр\\
RodoSol-ALPR & 20\,000 & хайрцаг + текст & Бразил/Mercosur, академик\\
Roboflow Vehicle Registration Plates (HF) & 8\,823 & хайрцаг & олон улс, нээлттэй\\
Roboflow mongolian-plate~\cite{roboflowmn} & 984 & хайрцаг & Монгол, CC BY 4.0, API түлхүүр\\
gumbernator/Mongolian-ALPR~\cite{gumbernator} & 68 & хайрцаг + тэмдэгт & Монгол, Apache-2.0\\
MLUB car-plate-prediction (via~\cite{tsund}) & 7\,394 & текст & Монгол, тэмцээн\\\bottomrule
\end{tabular}}\end{table}
